% Section S12, first subsection: the four results marked S in Sections 4 and 5 of the main text,
% with every definition they need and complete hand proofs. Results are numbered A.1, A.2, ... by hand.
\subsection*{The Four Main Results, Self-Contained}

This subsection states and proves the four results of the main text that carry the label S: Theorem~\ref{thm:validity},
Corollary~\ref{cor:v4}, the certificate statement of Section~\ref{sec:kings-valid} and Theorem~\ref{thm:hss}. It can be
read without the rest of this section. All definitions are given here, and every proof is a hand proof that uses only
these definitions and earlier steps of this subsection. No step rests on a computation. The one computed statement we
quote (part (c) of Proposition~A.14) is marked as such and is used nowhere. Numbered items are A.1, A.2, and so on.
We assume the basic theory of short games: disjunctive sums, outcome classes, and Zermelo's theorem that a finite game
without draws has a winner~\cite{Siegel2013}.

\subsubsection*{A.I Board, pawns and passed pawns}

\textbf{Board.} The board has $n$ files and $H\ge 3$ ranks. It is \emph{flat} (files $1,\dots,n$ in a row) or a
\emph{cylinder} (file $n$ is also next to file $1$). The \emph{distance} of two files is their distance in the row, or
on the cycle for the cylinder. White pawns move up, black pawns down. A square is \emph{ahead} of a white pawn if its rank is higher than the pawn's, and ahead of a black pawn if
its rank is lower. No pawn stands on its own first rank. In this subsection a position holds pawns and at most two
kings, and no other piece until a promotion.

\textbf{Pawn moves.} A pawn steps one square forward to an empty square; from its own second rank it may instead
\emph{double step} two squares forward over an empty square to an empty square; it captures an enemy pawn diagonally
forward on a file at distance $1$ (across the seam on the cylinder). If the last move was a double step of a pawn $q$ from
$a$ over $b$ to $c$, an enemy pawn on the square beside $c$ (same rank, file at distance $1$) may capture \emph{en passant}
(e.p.): it moves to $b$ and $q$ is removed. This is allowed on the next move only. The position then carries a
\emph{right}; the right is \emph{usable} if such an enemy pawn exists. A pawn that reaches its last rank promotes; in the pawn games $W$ and $S$
below it stays there as a pawn, as in our engine. Kings are
never captured.

\textbf{A.1 Lemma (files).} \emph{A pawn changes file only by a capture, and it then lands on a file that held an enemy
pawn just before the move. Hence, before the first promotion, no pawn ever enters a file that holds no pawn.}

\begin{proof}
A step or double step keeps the file. A capture lands on the square of the captured pawn; an e.p.\ capture lands on $b$,
on the file of the captured pawn $q$. Before the first promotion the only capturable pieces are pawns.
\end{proof}

\textbf{Components.} A \emph{component} of a pawn position is a maximal set of consecutive files that all hold pawns (an
arc on the cylinder), with the pawns on them; on a cylinder with no empty file the whole board is one component, a
\emph{ring}. The files next to a component are empty. In play with pawns only, a move of a pawn in a component $D$ uses only squares on the files
of $D$ (Lemma~A.1), so it changes no other component, and whether it is legal depends only on $D$ and on the right.

\textbf{Passed pawns.} A pawn passes the \emph{textbook test} if no enemy pawn stands ahead of it on its own file or on a
file at distance~$1$. It passes the \emph{e.p.-aware test} if it passes the textbook test and it is not the case that it
has just double stepped with an enemy pawn beside it (Definition~\ref{def:epaware}). A \emph{setting} fixes the test
(textbook, e.p.-aware, or e.p.\ off: no e.p.\ captures and the textbook test) and a \emph{race rule}, which names a winner
when both sides have a passed pawn. We only use that the race rule is \emph{local}: its verdict depends only on the passed
pawns, the pawns on their files and on the files at distance~$1$, and on which side made the last move. The count, mover and
opponent rules of Section~\ref{sec:rules} are local. Under every test, whether a pawn of a component is passed depends only
on that component and on the right, since the files at distance $1$ from its file belong to the component or are empty.

\textbf{A.2 Lemma.} \emph{(a) A pawn on its last rank is passed under every test. (b) Let $\Phi_t$ be the sum, over the
pawns of side $t$, of the number of single steps the pawn needs to reach its last rank. Every
pawn move of $t$ lowers $\Phi_t$ by at least $1$, and no move raises $\Phi_t$. (c) Let $H\ge5$. If a position has no
passed pawn and its right, if any, was created by the last move of the side not to move, the side to move has no
promoting move.}

\begin{proof}
(a) No square lies ahead of it, so it passes the textbook test. It fails the e.p.\ clause only if it has just double
stepped with an enemy pawn beside it; that pawn would stand on its own first rank, where no pawn stands. (b) Each pawn move
advances the pawn by one or two ranks; a capture removes a pawn of the other side. (c) Since $H\ge5$, a double step
does not reach the last rank, so a promoting pawn stands on the rank before its last. Only the last rank lies ahead of
it, and it holds no enemy pawn, so the pawn passes the textbook test. The last move was the opponent's, so the e.p.\ clause
does not apply, and the pawn is passed.
\end{proof}

\noindent Both restrictions in (c) are needed. On four ranks, in \pos{W a2 / B b3} no pawn is passed and a2--a4
promotes. On five ranks with the e.p.-aware test, after a2--a4 in \pos{W a2 b2 / B b4} no pawn is passed, yet a4--a5
promotes; White can play it only by moving twice in a row, which happens in a sum. Part (a) is what makes $W$ and
$S$ well defined in all these cases: a promoting move makes a passed pawn.

\subsubsection*{A.II The whole-board game, the sum model and the hypotheses of Theorem~\ref{thm:hss}}

We write White $=$ Left and Black $=$ Right, and $\bar\sigma$ for the opponent of a side $\sigma$. A positive integer is
the game $k=\game{k-1}{}$, with $0=\game{}{}$: Black has no move in it, and White has exactly $k$ moves in it before it is
$0$. We write $\sigma\BIG$ for $+\BIG$ if $\sigma$ is White and $-\BIG$ if $\sigma$ is Black.

\textbf{Whole-board game $W(X)$.} Fix a setting. Let $X$ be a pawn position, possibly with a right, that has no passed
pawn. A position of $W(X)$ is a pawn position with its right and the side to move. A move is a legal pawn move. If after a
move some side has a passed pawn, the game ends: if only $\sigma$ has one, $\sigma$ wins; if both have one, the race rule
names the winner. Otherwise play goes on, and a player with no move loses. By Lemma~A.2(b), $\Phi_W+\Phi_B$ falls at every
move, so $W(X)$ is finite and has no draws; by Lemma~A.2(a) a promoting move ends it. We write $\mathrm{Reach}(X)$ for the
positions of $W(X)$ reached from $X$ by play, with either side to move first, before the game ends ($X$ included). The
game $W(X)+\nim{h}$ is played in the same way, with one more move available: a move in the heap. A passed pawn ends the whole
game with the same winner, and a player with no move in $W(X)$ and an empty heap loses. (Replacing the end by an integer
$\pm\BIG$ with $\BIG>0$ and continuing in the heap gives the same outcomes, since $\BIG+\nim{h}$ is won by White whoever
starts: White empties the heap or moves in $\BIG$, and Black has no move in a nonnegative integer.)

\textbf{Sum model $S(X)$.} For a component $D$ (with its right, if the right concerns its pawns) let
\[
S(D)=\game{S(D^L_1),\,S(D^L_2),\dots}{S(D^R_1),\dots},
\]
where $D^L_i$ runs over the positions on the files of $D$ after a legal white move in $D$ (with the right that move
creates), $D^R_j$ likewise for black, and for such a position $D'$ we set $S(D')=\sigma\BIG$ if some side has a passed pawn
in $D'$ and $\sigma$ is the winner as in $W$, and otherwise $S(D')=\sum_E S(E)$ over the components $E$ of $D'$. Then
$S(X)=\sum_D S(D)$ over the components of $X$. The recursion ends by Lemma~A.2(b), and by Lemma~A.2(a) a promoting move
freezes its component. This is the main text's $S(X)$. Note the order: after a move the whole moved component is tested first, and
frozen at $\pm\BIG$ if a side has a passed pawn in it; only otherwise is it re-split. (Splitting first would be a different model: in \pos{W b2 d6 / B c7} on the ordinary board, c7$\times$d6 leaves two passed pawns on separated files, a race that White wins under the count rule, whereas $+\BIG-\BIG=0$.) In a sum the right of a component
stays with that component while the players move elsewhere; in $W(X)$ a right lapses after one move. This difference is
what condition (E) controls.

\textbf{Move bound.} $N(X)$ is the largest number of moves of one side in a finite sequence of legal pawn moves from $X$,
made by either side in any order, where e.p.\ is allowed only directly after the double step it answers and play goes on
past passed pawns. By Lemma~A.2(b), $N(X)\le\max(\Phi_W(X),\Phi_B(X))$, and a sequence from a position of
$\mathrm{Reach}(X)$, prefixed by the play that reached it, is a sequence from $X$.

\textbf{Threats, HSS and (E).} A component $D$ with no passed pawn \emph{carries a threat for $t$} if there is a finite
sequence of legal pawn moves on the files of $D$, starting from $D$ with its right, made by either side in any order, such
that no position before the last has a passed pawn and the last position gives $t$ a passed pawn (a race gives both sides
one). The right of each position in the sequence is the one created by the move before it.
$X$ is \emph{hereditarily sum-safe} (HSS) if no $Y\in\mathrm{Reach}(X)$ has two components that carry threats for opposite
sides. $X$ satisfies \emph{(E)} if no $Y\in\mathrm{Reach}(X)$ with at least two components carries a usable right.

\textbf{A.3 Remark (the precise form is implied by the main text's).} Definition~\ref{def:hss} and Theorem~\ref{thm:hss}
quantify over positions reachable by moves of either side \emph{in any order}; we use only $\mathrm{Reach}(X)$, which is
smaller, so our HSS and (E) are implied by the main text's. The main text's threat does not say ``before the first passed
pawn''. Reading it without that clause gives more threats and a stronger hypothesis. Our reading has fewer threats, so the
theorem below covers both readings. ``$\BIG$ exceeds the number of moves either side can make'' is $\BIG>N(X)$. The main
text leaves implicit that the race rule is local, and that a pawn on its last rank stays a pawn and so counts as passed, as in our engine.

\textbf{A.4 Lemma (what (E) gives).} \emph{Let $X$ satisfy (E), let $Y\in\mathrm{Reach}(X)$ have at least two components,
and let $\mu$ be a move in a component $C$ of $Y$, giving $Y'$. For every other component $D$ of $Y$: (a) the games
$S(D)$ with the right of $Y$ (if it concerns $D$) and without it have the same options; (b) no pawn of $D$ is passed in
$Y'$.}

\begin{proof}
(a) By (E) the right of $Y$ is not usable, so it allows no move, and every other move of $D$ leads to the same position
with the same new right. (b) Pawns of $D$ are not passed in $Y$, since the game did not end there. A pawn of $D$ that
passes the textbook test in $Y$ would have failed only the e.p.\ clause, so $Y$ would carry a usable right, against (E). So
no pawn of $D$ passes the textbook test in $Y$, and $D$ is unchanged in $Y'$.
\end{proof}

\textbf{A.5 Lemma (threats are inherited).} \emph{Let $E$ be a component with no passed pawn that carries no threat for
$\bar\sigma$. Let $\bar\sigma$ make a move $\mu$ in $E$, giving $E'$ on the files of $E$. Then $S(E')=\sigma\BIG$, or $E'$ has no passed pawn and each of its components carries no threat for
$\bar\sigma$.}

\begin{proof}
If $\bar\sigma$ has a passed pawn in $E'$ (alone or in a race), the one-move sequence $\mu$ shows that $E$ carries a
threat for $\bar\sigma$. So either only $\sigma$ has one, and $S(E')=\sigma\BIG$, or $E'$ has none. In the second case
suppose a component $E_i$ of $E'$ carries a threat for $\bar\sigma$, by a sequence $s$. We show that $\mu$ followed by $s$
witnesses a threat of $E$ for $\bar\sigma$. The moves of $s$ use only the files of $E_i$, so they are legal after $\mu$. If
$\mu$ is a double step, it changes no file's occupancy, so $E'$ has the files of $E$ and $E_i=E'$, with the right of $\mu$,
which is the right the sequence $s$ starts from. If $\mu$ is not a double step, $E'$ carries no right, and every other
component of $E'$ stays as it is in $E'$, with no passed pawn and no right, during $s$. In both cases every position before
the last has no passed pawn ($E$, $E'$, and those of $s$), and the last one gives $\bar\sigma$ a passed pawn. This
contradicts the hypothesis on $E$.
\end{proof}

\textbf{A.6 Lemma (frontier).} \emph{Let $X$ satisfy (E) and $\BIG>N(X)$. Let $Y\in\mathrm{Reach}(X)$, and let a move
$\mu$ in a component $C$ of $Y$ end $W$ with winner $\sigma$. Suppose no component $D\ne C$ of $Y$ carries a threat for
$\bar\sigma$. Then $\sigma$ wins $G=\sigma\BIG+\sum_{D\ne C}S(D)$ whoever moves first.}

\begin{proof}
Take $\sigma$ White; the other case is the mirror image. We play on the game trees as displayed, without replacing a pawn
component by a simpler game of the same value. White's strategy is to move, at each turn, in the original summand $+\BIG$, one move at a time.
Black has no move in a nonnegative integer, so every Black move is a pawn move in a summand $S(E)$, where $E$ is
a component of $Y$ other than $C$ or arises from one. By Lemma~A.5 and induction on the number of Black moves, each such $E$
has no passed pawn and carries no threat for Black, and every component that Black's move freezes becomes $+\BIG$. Black never captures
e.p.: White makes no pawn move after $\mu$, and the only right White may have created earlier is the right of $Y$, which
is not usable by (E), since $Y$ has the two components $C$ and $D$ (if $Y$ has no component $D\ne C$, Black has no pawn move
at all). So Black's moves, in order, follow the play from $X$ to $Y$ and $\mu$ as a sequence of legal pawn moves, and Black
makes at most $N(X)<\BIG$ of them. White needs at most one move per Black move, plus one if White starts, so at most
$N(X)+1\le\BIG$ moves, and the summand $\BIG$ alone supplies $\BIG$ of them. So White always has a move. If Black makes $k$ moves in all, White needs $k$ moves when Black starts and $k+1$ when White starts, and $k+1\le N(X)+1\le\BIG$. The play is finite
(the integers shrink and $\Phi_B$ falls), so it ends with a player who cannot move, and this is Black.
\end{proof}

\textbf{A.7 Theorem (sum safety; Theorem~\ref{thm:hss}).} \emph{Main text: ``Suppose that $X$ has no passed pawn and is HSS, that no position
reachable from $X$ with two or more components has a usable en passant right (condition (E)), and that $\BIG$ exceeds the
number of moves either side can make. Then $S(X)$ and $W(X)$ have the same outcome.'' Precise form: let $X$ be a pawn
position with no passed pawn, fix a setting with a local race rule, and suppose that $X$ is HSS, satisfies (E), and
$\BIG>N(X)$. Then for each side to move first, $S(X)$ and $W(X)$ have the same winner.}

\begin{proof}
For $Y\in\mathrm{Reach}(X)$ let $S_Y=\sum_D S(D)$ over the components of $Y$, the right of $Y$ going with the component
of the pawn that made the double step. Then $S_X=S(X)$. A position of $\mathrm{Reach}(X)$ other than $X$ comes with the side $t$ to move in the play that reached it; at $X$ either side may be to move. We show, by induction on $\Phi_W+\Phi_B$, that $W$ at $Y$ and
$S_Y$ have the same winner with $t$ to move. (The other side to move at $Y$ is not needed and need not satisfy HSS or (E).)

The moves of $S_Y$ are the moves in the components of $Y$. These are exactly the legal moves of $W$ at $Y$: pushes and
captures use the files of one component (Lemma~A.1), and an e.p.\ capture uses the right of $Y$, which lies in the
component of the pawn that double stepped, next to the capturer. Fix a move $\mu$ in $C$, giving $Y'$ in $W$ and $C'$ on
the files of $C$. The option of $S_Y$ is $\sum_{D\ne C}S(D)+S(C')$.

\emph{$Y'$ has no passed pawn.} Then $C'$ has none, and $S(C')$ is the sum over the components of $C'$. The components of
$Y'$ are those of $C'$ and the $D\ne C$. If $Y$ has two or more components, Lemma~A.4(a) lets us drop the right of $Y$
from the summands $S(D)$ without changing them. So the option is $S_{Y'}$, and the induction hypothesis applies.

\emph{$Y'$ ends $W$ with winner $\sigma$.} In $W$, $\sigma$ wins at once. By Lemma~A.4(b) (or because $C$ is the only
component) every passed pawn of $Y'$ lies in $C'$, so $C'$ has a passed pawn, and $S(C')=\sigma\BIG$ because the race rule
is local and the mover is the same. The move $\mu$ alone shows that $C$ carries a threat for $\sigma$ (and for both sides
in a race). By HSS no other component of $Y$ carries a threat for $\bar\sigma$, and Lemma~A.6 shows that $\sigma$ wins the
option whoever moves first.

So every option has the same status in both games: won by the same player when the opponent moves next. A player with no
option loses in both. Hence the winners agree at $Y$, and at $Y=X$ the claim is the theorem.
\end{proof}

\textbf{A.8 Proposition (opposite threats cancel).} \emph{For every integer $\BIG\ge1$,
$\game{\BIG}{0}+\game{0}{-\BIG}=0$, although in the game where a $\pm\BIG$ ends play the first player wins.}

\begin{proof}
If White starts and moves to $\BIG+\game{0}{-\BIG}$, Black answers $-\BIG$, and White faces $0$; if White moves to
$\game{\BIG}{0}+0$, Black answers $0$. Either way White loses. Black starting loses in the same way. A second-player win
equals $0$. If instead reaching $\pm\BIG$ ends play, the first player moves to its own $\pm\BIG$ and wins.
\end{proof}

\textbf{A.9 Remark (the move bound).} A file contributes at most $1+\dots+(H-2)=(H-1)(H-2)/2$ to $\Phi_t$, when its squares on ranks $2$ to $H-1$ all hold pawns of $t$, so
$N(X)\le\max_t\Phi_t(X)\le n(H-1)(H-2)/2$, which is at most $168<256$ on boards with $n\le8$ files and $H\le8$ ranks. So $\BIG=256$ meets
$\BIG>N(X)$ there, as the main text says.

\subsubsection*{A.III King groups, the promotion game and the validity theorem}

\textbf{King group.} A \emph{king group} $Z$ is a set $F_Z$ of consecutive files (an arc on the cylinder) holding both
kings and some pawns, the \emph{Z-pawns}. The files at distance $1$ from $F_Z$ form the \emph{buffer} and hold no pawn.
The other pawns, the \emph{context}, stand on files at distance at least $2$ from $F_Z$. A \emph{Z-move} is a move of a
king or of a Z-pawn. $X$ is the pawn position of the context alone; we assume it has no passed pawn. In this part $H\ge5$, so that Lemma~A.2(c) applies.

\textbf{Three games.} \emph{Chess} $\mathrm{C}$ is played with the usual rules, any promotion, no fifty-move rule;
checkmate wins and stalemate and infinite play are draws. A side \emph{wins} a position if it has a strategy that ends every
play in a win for it in finitely many moves. The \emph{promotion game} $\mathrm{P}$ is chess in which the first promotion
ends the game and wins for the promoter. The \emph{frozen-king game} $\mathrm{F}$ is $\mathrm{P}$ in which, in addition,
a king move loses at once for its mover; our solver uses it to play the convention, and in it every king move is a loss by
definition.

\textbf{A.10 Lemma (buffer).} \emph{Before the first promotion, in any play from a position with a king group, no pawn
stands on a buffer file, Z-pawns stay on $F_Z$, context pawns stay on files at distance at least $2$ from $F_Z$, no context
pawn attacks a square of $F_Z$, and a Z-pawn captures only on $F_Z$.}

\begin{proof}
By Lemma~A.1 a pawn enters a file only if an enemy pawn stands on it. The buffer starts empty, so by induction on the moves
it stays empty. To leave $F_Z$, or to reach $F_Z$ from the context, a pawn would have to enter a buffer file. A pawn
attacks only the files at distance $1$ from its own, and a Z-pawn's capture lands on the file of the captured pawn (also in an e.p.\ capture), which is not a buffer file.
\end{proof}

\textbf{Decisive events.} Let $Q$ be a position made of a king group $Z$ and its context, with a side to move. A move is a
\emph{decisive event} if it is a Z-move, or a context move after which $W(X)$ would end, that is, some side has a passed
pawn among the context pawns. Since the Z-pawns are at least two files away, this is also the passed-pawn test on the whole
board for the context pawns. $\mathrm{R}(Q)$ is the set of chess positions reached from $Q$ by legal moves none of which is
a decisive event ($Q$ included). Kings and Z-pawns do not move along such play, so a position $Y\in\mathrm{R}(Q)$ is fixed
by its context part $X_Y$, its right and its side to move.

\textbf{Conditions.} At every $Y\in\mathrm{R}(Q)$, with side $s$ to move:
\begin{itemize}
\item[(V1)] the legal chess moves of $s$ that are not Z-moves are exactly the moves of $s$ in $W$ at $X_Y$, with the same
resulting context and right; none of these moves gives check and nobody is in check; and $s$ has a legal chess move (a Z-move counts);
\item[(V2)] every context move at $Y$ after which $W$ ends with winner $\sigma$ leads to a chess position that $\sigma$ wins;
\item[(V3)] every legal Z-move at $Y$ leads to a chess position that the opponent of its mover wins.
\end{itemize}
(V4) is: $X$ is HSS, satisfies (E), and $\BIG>N(X)$.

\textbf{A.11 Theorem (validity of the convention; Theorem~\ref{thm:validity}).} \emph{Main text: under (V1)--(V3)
``the chess result, with either side to move, is the outcome of $W(X)$ \dots\ without draws''. Precise form: let $Q$ be a
king group with its context, with no passed pawn in $X$, and let (V1)--(V3) hold at every position of $\mathrm{R}(Q)$. If
$\sigma$ wins $W(X)$ with the side to move of $Q$, then $\sigma$ wins $Q$ in chess. In particular $Q$ is not a draw.}

\begin{proof}
$W(X)$ is finite without draws, so some $\sigma$ wins it; fix a winning strategy. We give $\sigma$ a chess strategy and
show that every play under it is a $\sigma$-win. Along the play, as long as no decisive event has occurred, the position
$Y$ lies in $\mathrm{R}(Q)$ and $X_Y$ is a position of $W(X)$ won by $\sigma$ with the same side to move. This holds at $Q$.

\emph{$\sigma$ to move.} $\sigma$ plays the move of its strategy at $X_Y$. By (V1) it is a legal chess move with the same
result. If it ends $W$, the winner is $\sigma$, and by (V2) the chess position reached is won by $\sigma$; $\sigma$ follows
a winning strategy from there. Otherwise the move is not decisive, the new position is in $\mathrm{R}(Q)$, and $W$ is won by
$\sigma$ there.

\emph{$\bar\sigma$ to move.} By (V1), $\bar\sigma$ has a legal move. A Z-move leads to a position won by $\sigma$, by (V3).
Any other legal move is, by (V1), a move of $W$ at $X_Y$, and all of these lead to positions won by $\sigma$ in $W$. If it
ends $W$, the winner is $\sigma$, and (V2) applies. Otherwise the new position is in $\mathrm{R}(Q)$ and won by $\sigma$ in
$W$. (If $\bar\sigma$ has no move in $W$, its legal moves are all Z-moves.)

Every move before the first decisive event is a pawn move of $W$, so by Lemma~A.2(b) there are at most
$\Phi_W(X)+\Phi_B(X)$ of them; (V1) then still supplies a legal move, so a decisive event occurs by move $\Phi_W(X)+\Phi_B(X)+1$, and at the event (V2) or (V3) gives $\sigma$ a strategy that wins in finitely many further moves. The play is never stalemate (V1), and no promotion occurs
before the event (Lemma~A.2(c)).
\end{proof}

\textbf{A.12 Remark.} (a) The proof uses from (V1) only the equality of move sets and the existence of a legal move, not the
clauses on checks. A context move never gives check before a decisive event, by Lemma~A.10, since the kings stay on $F_Z$;
and a side in check from a Z-pawn has no legal context move, so the equality of move sets then fails unless $W$ has no move
either. (b) The theorem and the proof hold word for word
with chess replaced by $\mathrm{P}$ or by $\mathrm{F}$; in $\mathrm{F}$, (V3) holds for every king move by definition.
(c) $X$ must have no passed pawn; otherwise $W(X)$ is not defined.
(d) Unlike (V1a) of Theorem~1.14 in the longer treatment, (V1) contains no condition (E): it compares chess with the whole-board game $W(X)$, whose e.p.\ rights lapse after one move as in chess. Condition (E) is needed only to pass from $W(X)$ to $S(X)$ (Theorem~A.7), hence in (V4).

\textbf{A.13 Corollary (Corollary~\ref{cor:v4}).} \emph{Main text: ``If also (V4) \dots\ the chess result is the outcome
class of the components' $\pm\BIG$ value.'' Precise form: for a side to move for which the hypotheses of Theorem~A.11 hold, assume also (V4) and a local race rule, as in Theorem~A.7. Then the chess winner of $Q$ with that side to move is the winner of $S(X)$ with that side moving first. (The main text's ``either side to move'' assumes the hypotheses for both. They can hold for one side only: on the ordinary board, with the king group \pos{W Kb6 c7 / B Ka8 a5} and no other pawn, they hold with Black to move, who must play a5--a4 and is mated by c8=Q, but not with White to move, who mates at once although the empty sum is lost for the side to move.)}

\begin{proof}
By Theorem~A.11 the chess winner is the winner of $W(X)$, and by Theorem~A.7 this is the winner of $S(X)$.
\end{proof}

\textbf{A.14 Proposition (Proposition~\ref{prop:sumcounter}).} \emph{On $\cyl{8}$ with e.p.\ off, let
$X=$\,\pos{W b2 d4 g3 h5 / B a7 c7 f7 g4}. (a) $X$ is one component. (b) $X$ is not HSS. (c) (Checked by computer only, not
proved here and not used here.) $S(X)$ is a black win whoever starts and $W(X)$ a first-player win.}

\begin{proof}
(a) Only file e is empty. (b) White h5--h6 and Black a7$\times$h6 (across the seam) give
$Y=$\,\pos{W b2 d4 g3 / B c7 f7 g4 h6}. No pawn in $X$, after h5--h6, or in $Y$ is passed: g3 and g4 block each other, f7
has g3 ahead on a neighbouring file, c7 has b2 ahead, b2 and d4 have c7 ahead, a7 has b2 ahead, the white h-pawn has a7
ahead, and the black pawn on h6 has g3 ahead.
Files a and e of $Y$ are empty, so $Y$ has the components $\{f,g,h\}$ and $\{b,c,d\}$. The first carries a threat for
Black: f7--f6, f6--f5, f5--f4 keep g3 ahead of the f-pawn, and f4$\times$g3 leaves no white pawn in the component, so the new g3 pawn is passed. The second carries
one for White: d4--d5, d5--d6 keep c7 ahead of the d-pawn, and c7$\times$d6 empties file c, so b2 is passed (d6 is passed too, a race, which is a threat for both sides).
\end{proof}

\subsubsection*{A.IV The certificate}

\textbf{Data.} Let $Y$ be a position with a king group $Z$ and its context, no promotion yet, and side $s$ to move; write
$\bar s$ for the other side. Let $m$ be a legal Z-move of $s$ that does not promote (it may take $s$'s king off $F_Z$), and
$r$ a legal reply of $\bar s$ that is a Z-move, does not promote, and leaves $\bar s$'s king on $F_Z$. Let $Q$ be the
position after $m$ and $r$, with $s$ to move, and let $h\ge1$. The \emph{projection} $\pi(R)$ of a position $R$ deletes
every pawn off $F_Z$ and keeps the kings, the Z-pawns, the side to move, and a right whose captured pawn is on $F_Z$.

\textbf{Certificate of horizon $h$.} The move $m$ has one, with reply $r$, if (C1)--(C3) hold.
\begin{itemize}
\item[(C1)] \emph{Local promotion.} In the \emph{local game} from $\pi(Q)$, with $s$ to move, played by the rules of chess
except as follows, $\bar s$ has a strategy that wins against every play of $s$: $s$ may pass whenever it is not in check
(a pass changes only the side to move and clears the right); $\bar s$ may not pass and its king may only step onto $F_Z$;
the first promotion ends the game, a win for $\bar s$ if $\bar s$ promotes and a failure for $\bar s$ if $s$ does; a
checkmate of $s$ is a win for $\bar s$; $\bar s$ having no permitted move is a failure; and $\bar s$ must win with one of
its first $h$ moves, or it fails.
\item[(C2)] \emph{Counter-promotion.} Let $Q^\circ$ be $Q$ with $\bar s$'s king and all Z-pawns deleted. In $Q^\circ$
only $s$ moves: its king steps to squares that hold no pawn of $s$ and are not attacked by a pawn of $\bar s$ (it may
capture such pawns), and its pawns move by the pawn rules without e.p. Every sequence of such moves that ends with a
promotion has at least $h+1$ moves.
\item[(C3)] \emph{No stalemate.} $\bar s$ has a \emph{guided} strategy from $Q$, one that at each turn makes a Z-move
whose projection is a winning move of the local game of (C1) at the projected position, with the moves that remain, such that in no play from $Q$
following it is $s$ stalemated before $\bar s$'s $h$-th move.
\end{itemize}

\textbf{A.15 Remark (the main text's form).} The main text says that $m$ has a certificate of horizon $h$ ``if, with
only $Z$ on the board, the punisher promotes or mates within $h$ moves of its reply, no defender pawn outside $Z$ is that
fast, and no stalemate comes first''. This is a summary of (C1)--(C3): ``with only $Z$ on the board'' is the local game of
(C1), in which the context is deleted, the defender may pass and the punisher's king stays on $F_Z$; ``no defender pawn
outside $Z$'' is (C2), the defender's context pawns moving alone in $Q^\circ$ (its Z-pawns are covered by (C1)); and the
stalemate clause is (C3). A clause that an unsafe queen is a failure, used in the long version with chess in view, only strengthens
(C1); it gives no chess win either (Remark~A.17).

\textbf{A.16 Theorem (certificate).} \emph{Main text: ``In the promotion game (the first promotion wins) such a move loses
for its mover.'' Precise form: if $m$ has a certificate of horizon $h$ with reply $r$, then after $m$ the side $\bar s$ wins
$\mathrm{P}$: it plays $r$ and the guided strategy of (C3), and promotes or mates with one of its next $h$ moves.}

\begin{proof}
Consider any play from $Q$ in which $\bar s$ follows the guided strategy $\rho$ and $s$ makes legal moves, up to $\bar s$'s
$h$-th move or the first promotion. Let $R$ denote its positions.

\emph{Invariants.} (I1) By Lemma~A.10 (from $Y$; the kings' squares play no role in it), the buffer has no pawn, Z-pawns
stay on $F_Z$, context pawns stay off $F_Z$ and the buffer, no context pawn attacks a square of $F_Z$, and Z-pawns capture
only on $F_Z$. (I2) $\bar s$ makes only Z-moves and its king stays on $F_Z$. (I3) $\bar s$'s context pawns never move and
are only removed. No $\bar s$ pawn of the context attacks $s$'s king: none does after $m$, since $m$ is legal; $r$ moves no
context pawn; later the set of squares these pawns attack only shrinks, and $s$'s king never steps onto an attacked square.
So $s$ is in check in $R$ exactly when a Z-pawn of $\bar s$ attacks its king, which is the same in $\pi(R)$; the kings
never attack each other. (I4) Rights: a right created by $\bar s$ comes from a Z-pawn and its captured pawn is on $F_Z$; a
right created by a context double step of $s$ could be used only by a $\bar s$ context pawn, which never moves. So the only
e.p.\ captures in the play are by Z-pawns of pawns on $F_Z$, and $\pi$ keeps those rights.

\emph{Step 1: moves of $s$ project to local moves or passes.} Let $s$ be to move in $R$. A king move of $s$ is legal in
$\pi(R)$: the target square holds the same piece or less, the kings are the same, and the pawns of $\bar s$ that attack it
in $\pi(R)$ also do in $R$; a capture of a context pawn becomes a plain move. A Z-pawn move of $s$ uses squares of $F_Z$,
which carry the same pieces in $R$ and $\pi(R)$, and the checks are the same by (I3); a Z-pawn promotion is a promotion in
the local game. A context-pawn move leaves $\pi(R)$ unchanged except for the side to move and the right, as a pass does; it
is legal only if $s$ is not in check, since a context pawn can neither capture on $F_Z$ (I1) nor block a pawn check; so the
pass is allowed in $\pi(R)$. Hence the projection of the play is a play of the local game in which $\bar s$ plays only
winning moves, with the same count of moves of each side.

\emph{Step 2: permitted local moves of $\bar s$ are legal.} Let $\bar s$ be to move in $R$ and let $\nu$
be a permitted local move at $\pi(R)$. Its squares lie on $F_Z$, which holds the same pieces in $R$ and $\pi(R)$, and its
right is the same (I4). After $\nu$, $\bar s$'s king stands on $F_Z$, where no context pawn of $s$ attacks (I1); the other
attackers are the same as in $\pi(R)$. There are no sliding pieces, so there are no pins. So $\nu$ is legal in $R$. By Step~1 the projected position is one that $\bar s$ wins in the local game with the moves that remain, so a winning move exists there, and the
guided strategy is defined at every turn of $\bar s$ before it wins; in particular $\bar s$ is never mated or stalemated.

\emph{Step 3: $s$ does not promote first.} By Step 1 and since $\bar s$, playing only winning moves, wins the local game by its $h$-th move, which
follows $s$'s $h$-th turn, $s$ has at most $h$ turns. A promotion of a Z-pawn of $s$ would be a promotion by $s$ in the
local game, a failure for $\bar s$, which cannot happen from a local position that $\bar s$ wins. Suppose a context pawn of $s$ promotes at
$s$'s $j$-th turn, $j\le h$. Delete from the play all moves of $\bar s$ and all Z-pawn moves of $s$. What remains is at most
$j$ moves of $s$, king moves and context-pawn moves, ending with the promotion. It is a legal sequence in $Q^\circ$: $s$'s
king and the context pawns stand on the same squares in both, since only $s$ moves them; a king move in the play goes to a
square not attacked by a context pawn of $\bar s$ and not holding a pawn of $s$; a context-pawn move uses context squares,
whose occupants are the same, since $\bar s$'s king and the Z-pawns are on $F_Z$; and no e.p.\ capture is involved (I4).
This contradicts (C2).

\emph{Step 4: $\bar s$ wins.} $s$ cannot mate $\bar s$ (Step 2), $\bar s$ is not stalemated (Step 2), $s$ is not
stalemated before $\bar s$ wins (C3), and $s$ does not promote first (Step 3). The projected play is a play of the local
game in which $\bar s$ plays only winning moves, so it ends by $\bar s$'s $h$-th move with a promotion by $\bar s$ or a mate of $s$. A
promotion is a legal real move (Step 2) and wins $\mathrm{P}$. If $s$ is mated in $\pi(R)$, it is in check in $R$ (I3),
and a legal move of $s$ in $R$ would project to a legal local move, a pass being impossible in check (Step 1); so $s$ is
mated in $R$. Every play is finite, since it lasts at most $h$ moves of $\bar s$.
\end{proof}

\textbf{A.17 Remark.} Theorem~A.16 is a statement about $\mathrm{P}$ only. A won promotion in $\mathrm{P}$ need not be a
won chess game, so it does not give (V3) for chess, as the main text says. With Remark~A.12(b) and $Q$ as in Theorem~A.11, certificates for every
Z-move at every position of $\mathrm{R}(Q)$ give (V3) in $\mathrm{P}$.

\textbf{Where to find the longer treatment.} In the numbering of the rest of this section: Lemma~A.1 and A.10 are
Lemma~2.2; Lemma~A.2(b) and the move bound are the tempo counts of Definition~1.10; Lemma~A.6 and Theorem~A.7 are the frontier step
and the induction of Theorem~1.11(c); Lemma~A.5 is the inheritance of threats by re-split components in that frontier step, and Lemma~A.4
is the use that the induction makes of (E);
Proposition~A.8 is Proposition~1.12; Theorem~A.11 is the chess form of Theorem~1.14(i), proved as the sharp form of Theorem~1.14 with $N=\mathrm{R}(Q)$ (since (V2) and (V3) are stated in chess, (A$_Q$) is not needed), Remark~A.12(a) is
Proposition~1.15(a), and Corollary~A.13 is the chess form of Theorem~1.14(ii); the certificate (C1)--(C3) is Definition~2.6, conditions (L1$_P$),
(L2), (L3), and Theorem~A.16 is Theorem~2.7(a); Remark~A.17 is ``What Theorem 2.7 does not give''.
